\documentclass[a4paper,11pt,oneside]{book}

%========================================================================%
\usepackage{courier}
\usepackage[utf8]{inputenc}   
\usepackage[T1]{fontenc}      

\usepackage{graphicx}
\usepackage{xstring}
\usepackage{etoolbox}
\usepackage{float}
\usepackage{tabularx}
\usepackage[font=small,skip=2pt]{caption}
\usepackage{subcaption}
\captionsetup{within=none}
\usepackage{wrapfig}
\usepackage{multicol}
\usepackage{eurosym}
\usepackage{footmisc}
\usepackage{hyperref}
\usepackage{hyperxmp}
\usepackage{booktabs}
\usepackage[toc,page,title,titletoc]{appendix} % úsalo solo si agregas apéndices
\renewcommand*\appendixpagename{Apéndices}

\usepackage[spanish]{babel}
\addto\captionsspanish{\renewcommand{\tablename}{Tabla}}
\addto\captionsspanish{\renewcommand{\listtablename}{Índice de Tablas}}

% TikZ se carga DESPUÉS de babel: el español vuelve activos los símbolos < y >,
% y la librería 'babel' de TikZ es la que evita que eso rompa las flechas.
\usepackage{tikz}
\usetikzlibrary{arrows.meta,babel}

\usepackage[dvipsnames]{xcolor}
\definecolor{azul}{RGB}{81, 131, 187}   % Azul institucional UNIR
\def\azul{\color{azul}}
\def\an{\bfseries\azul}

\usepackage{geometry}
\geometry{
    paper=a4paper,
    left=3.5cm,
    right=1.5cm,
    top=2.5cm,
    bottom=2.5cm,
    headheight=0.75cm,
    headsep=1.5cm,
    footskip=1.5cm,
}

\usepackage{lastpage}
\usepackage{fancyhdr}
\pagestyle{fancy}
\patchcmd{\chapter}{\thispagestyle{plain}}{\thispagestyle{fancy}}{}{}

\usepackage{array}
\usepackage{setspace}
\setstretch{1.5}
\usepackage[autostyle=true]{csquotes}

% Bibliografía en formato APA (7.ª ed.) real, vía el paquete biblatex-apa:
% sangría francesa, orden alfabético autor-año, "et al." según reglas APA.
\usepackage[backend=biber,style=apa,sorting=nyt]{biblatex}
\addbibresource{referencias.bib}
\renewcommand*{\cite}[1]{\parencite{#1}}
\newcommand*{\apa}[1]{\citeauthor{#1} (\citeyear{#1})}
\usepackage{ifxetex,ifluatex}
\newif\ifxetexorluatex
\ifxetex\xetexorluatextrue\fi
\ifluatex\xetexorluatextrue\fi

\ifxetexorluatex
    \usepackage{fontspec}
    \newfontfamily\georgiafont{TeX Gyre Pagella}   % sustituto libre de Georgia
    \newfontfamily\cambriafont{TeX Gyre Termes}    % sustituto libre de Cambria
    \newcommand{\georgiab}{\georgiafont\bfseries}
    \newcommand{\georgiamd}{\georgiafont}
    \newcommand{\cambriamd}{\cambriafont}
\else
    \newcommand{\georgiab}{\rmfamily\bfseries}
    \newcommand{\georgiamd}{\rmfamily\mdseries}
    \newcommand{\cambriamd}{\rmfamily\mdseries}
\fi

%========================================================================%
\makeatletter
\def\autor#1{\author{#1}}
\def\titulo#1{\title{#1}}
\def\fecha#1{\date{#1}}
\def\universidad#1{\gdef\@universidad{#1}}
\def\escuela#1{\gdef\@escuela{#1}}
\def\master#1{\gdef\@master{#1}}
\def\tipoTesis#1{\gdef\@tipoTesis{#1}}
\def\director#1{\gdef\@director{#1}}
\def\ciudad#1{\gdef\@ciudad{#1}}
\gdef\keywords#1{\gdef\keywordsv{#1}}
\gdef\keywordsEs#1{\gdef\keywordsESv{#1}}
\makeatother

\def\n{\bfseries}
\def\glmtr{\guillemotright}
\def\glmtl{\guillemotleft}

\newcolumntype{L}[1]{>{\hsize=#1\hsize\raggedright\arraybackslash}X}%
\newcolumntype{R}[1]{>{\hsize=#1\hsize\raggedleft\arraybackslash}X}%
\newcolumntype{C}[1]{>{\hsize=#1\hsize\centering\arraybackslash}X}%
\newcommand{\cc}[1]{\begin{minipage}[c][0.5cm][c]{0.5cm}\centering#1\end{minipage}}

\newenvironment{codigo}{
    \begin{minipage}[t][][t]{0.7\textwidth}
    \ttfamily \small \linespread{0}
}{
    \end{minipage}
}
\newcommand{\cod}{\ttfamily \small \linespread{0}}

\newenvironment{abstract}[1][Resumen]{%
    \begin{center}%
        {\fontsize{14}{14}\selectfont\bfseries #1\vspace{-0.5em}}%
    \end{center}%
    \quotation
}{\endquotation\vspace{1cm}}

%========================================================================%
\hypersetup{
    pdfpagemode={UseOutlines},
    bookmarksopen=true,
    bookmarksopenlevel=0,
    hypertexnames=false,
    colorlinks=true,
    pdfborder={0 0 0},
    citecolor=blue,
    linkcolor=blue,
    urlcolor=blue,
    pdfstartview={FitV},
    breaklinks=true
}

% Encabezado
\fancypagestyle{mainmatter}{%
    \fancyhead{}%
    \lhead{{\color{azul}Actividad grupal}}%
    \chead{}%
    \rhead{\masterv}%
    \lfoot{{\color{azul}Computación neuromórfica}}%
    \cfoot{}%
    \rfoot{\thepage\ de \getpagerefnumber{LastPage}}}

\fancypagestyle{frontmatter}{%
    \fancyhead{}%
    \lhead{{\color{azul}Actividad grupal}}%
    \chead{}%
    \rhead{\masterv}%
    \lfoot{{\color{azul}Computación neuromórfica}}%
    \cfoot{}%
    \rfoot{\thepage}}

%========================================================================%
\renewcommand*{\maketitle}{%
    \begin{titlepage}%
        {\flushleft%
            {\includegraphics[width=8.18cm]{logo-UNIR.pdf}}\par
        }%
        \vspace{1cm}
        \begin{table}[h!]
        \hyphenpenalty=10000
        \setlength\arrayrulewidth{2pt}
            \begin{tabularx}{14.5cm}{ L{0.074} !{\color{azul}\vline}  L{0.931} }
                & \\
                & \fontsize{14}{14}\selectfont \georgiab \universidadv\\
                & \\
                & \fontsize{18}{18}\selectfont \georgiab \escuelav \\
                & \\
                & \fontsize{14}{14}\selectfont  \georgiab\masterv \\
                & \\ & \\ & \\
                & \fontsize{42}{42}\selectfont \cambriamd \color{azul}   \titulov \\
            \end{tabularx}
        \end{table}
        \vspace{0.5cm}
        {\flushleft \fontsize{11}{11}\selectfont %
            \georgiab \tipoTesisv\par
                \vspace{0.5cm}
                Presentado por: \georgiamd \autorv \par
                \vspace{0.5cm}
                \georgiab Director: \georgiamd \directorv \par
                \vspace{3cm}
                {\color{azul}%
                \georgiab Ciudad: \georgiamd \ciudadv \par
                \georgiab Fecha: \georgiamd \fechav \par}
        }
    \end{titlepage}}

%========================================================================%
% NOTAS AL EQUIPO. Se ven en rojo en el PDF para no perderlas de vista.
% Al resolver cada una, bórrala. Antes de entregar no debe quedar ninguna.
\newcommand{\equipo}[1]{{\color{red}\textbf{[EQUIPO]} \textit{#1}}}

%========================================================================%
\titulo{¿Qué retos plantean los nuevos modelos de computación?}
\autor{Benjamín David Salín Monroy, Héctor Lobato Silva, Nora Arroyo Carranza, Josué Isaac Gómez Méndez}
\director{Dr. Miguel Alejandro Medina Armendariz}
\fecha{12-10-2026}
\ciudad{CDMX}
\universidad{Universidad Internacional de la Rioja (UNIR)}
\escuela{Escuela Superior de Ingeniería y Tecnología}
\master{Máster en Computación Cuántica}
\tipoTesis{Actividad grupal · Computación neuromórfica}

\keywords{Neuromorphic computing, spiking neural networks, technology ethics}
\keywordsEs{Computación neuromórfica, redes de impulsos, ética de la tecnología}


%========================================================================%
\makeatletter
    \xdef\autorv{\@author}
    \xdef\titulov{\@title}
    \xdef\masterv{\@master}
    \xdef\fechav{\@date}
    \xdef\universidadv{\@universidad}
    \xdef\escuelav{\@escuela}
    \xdef\tipoTesisv{\@tipoTesis}
    \xdef\directorv{\@director}
    \xdef\ciudadv{\@ciudad}
\makeatother
\hypersetup{
    pdftitle={\titulov},
    pdfauthor={\autorv},
    pdfsubject={\masterv},
    pdfkeywords={\keywordsESv},
}

%========================================================================%
%                        INICIO DEL DOCUMENTO
%========================================================================%
\begin{document}

    %--- Páginas iniciales: portada, índice, resumen, listas ---%
    \frontmatter
    \pagestyle{frontmatter}
        \maketitle

        \begin{abstract}[Resumen]
            \equipo{Redactar el resumen al final, cuando el resto esté cerrado:
            150--200 palabras que digan qué paradigma se analiza, qué desafíos
            y dilemas se identifican y a qué conclusión llega el equipo.}
        \end{abstract}

        \clearpage
        % Los índices van con interlineado sencillo para que quepan en una
        % sola hoja; el cuerpo del documento conserva el 1,5 de la plantilla.
        \begin{spacing}{1}
            \small
            \tableofcontents
            % \listoffigures abriría capítulo y forzaría página nueva; se
            % compone a mano para que el índice de figuras quede en la misma
            % hoja que el índice general.
            \vspace{1.6em}
            {\large\bfseries\color{azul} Índice de figuras}\par\vspace{0.7em}
            \makeatletter\@starttoc{lof}\makeatother
        \end{spacing}



    %========================================================================%
    %  CUERPO DEL DOCUMENTO — ESTRUCTURA OBLIGATORIA UNIR
    %  (introducción, desarrollo, resultados, discusión, conclusiones,
    %   bibliografía). Escriba el contenido de cada sección directamente
    %  aquí abajo; ya no hace falta ningún \include ni carpeta "contenido/".
    %========================================================================%
    \mainmatter
    \pagestyle{mainmatter}

\chapter{Introducción}

Durante más de siete décadas, la práctica totalidad de los sistemas de cómputo se
ha construido sobre la arquitectura de von Neumann, que separa la unidad de
procesamiento de la memoria. Esa separación obliga a transferir datos de forma
continua entre ambas, y ese tránsito, y no la velocidad de cálculo, se ha
convertido en el principal límite de rendimiento y consumo. El problema se agrava
con las cargas de trabajo de la inteligencia artificial, que mueven volúmenes de
parámetros muy superiores a los de la computación convencional.

La computación neuromórfica propone una alternativa estructural: diseñar hardware
inspirado en la organización del sistema nervioso, donde el procesamiento y el
almacenamiento coinciden físicamente y la comunicación se produce mediante
impulsos discretos, solo cuando hay actividad \parencite{mead1990}. No se trata
de ejecutar más rápido el mismo modelo de cómputo, sino de cambiar la forma de
organizarlo.

\begin{figure}[h!]
\centering
\begin{tikzpicture}[font=\small,
    panel/.style={draw=gray!55,rounded corners=3pt,
                  minimum width=5.6cm,minimum height=2.6cm},
    caja/.style={draw=gray,fill=white,rounded corners=2pt,
                 minimum width=1.9cm,minimum height=1.0cm,align=center},
    celda/.style={draw=azul,fill=azul!12,
                  minimum width=1.85cm,minimum height=1.0cm,align=center}]
  % --- izquierda: arquitectura convencional
  \node[panel] at (0,0) {};
  \node[anchor=south] at (0,1.45) {\textbf{Arquitectura convencional}};
  \node[caja] (cpu) at (-1.45,0.25) {Procesador};
  \node[caja] (mem) at ( 1.45,0.25) {Memoria};
  \draw[<->,line width=1.6pt,red!65!black] (cpu) -- (mem);
  \node[color=red!65!black] at (0,-0.72) {el dato viaja en \emph{cada} operación};
  % --- derecha: arquitectura neuromórfica
  \node[panel] at (7.4,0) {};
  \node[anchor=south] at (7.4,1.45) {\textbf{Arquitectura neuromórfica}};
  \node[celda] at (6.48,0.25) {Cómputo};
  \node[celda,fill=azul!25] at (8.32,0.25) {Memoria};
  \draw[azul,line width=1.1pt,rounded corners=2pt]
        (5.52,-0.33) rectangle (9.28,0.83);
  \node[color=azul] at (7.4,-0.72) {el dato ya está donde se opera};
\end{tikzpicture}
\caption{Separación entre procesamiento y memoria frente a su colocalización.}
\label{fig:arquitectura}
\end{figure}

Su relevancia actual responde a dos presiones convergentes. La primera es
energética: el coste de entrenar y ejecutar modelos de inteligencia artificial
crece a un ritmo que la infraestructura convencional difícilmente sostiene. La
segunda es la necesidad de procesamiento autónomo en el propio dispositivo, sin
depender de una conexión permanente con centros de datos remotos, en ámbitos como
la robótica, los vehículos autónomos o las prótesis médicas
\parencite{roy2019}.

\section{Objetivo y estructura de la memoria}

El objetivo de este trabajo es analizar los retos que plantea la computación
neuromórfica en tres planos: el tecnológico, el ético y el social.

\equipo{Añadir una cifra que ilustre el coste de mover un dato frente a
operarlo, y un dato verificado del crecimiento del consumo de los centros de
datos de IA. Ambos refuerzan la justificación del trabajo.}

La memoria se organiza como sigue. El capítulo \ref{cap:desarrollo} describe el
paradigma, sus fundamentos teóricos y sus aplicaciones prácticas. El capítulo
\ref{cap:resultados} expone los tres desafíos tecnológicos y las tres
implicaciones éticas identificadas. El capítulo \ref{cap:discusion} propone
medidas concretas para abordar esos riesgos y valora el impacto del paradigma en
la sociedad. El capítulo \ref{cap:conclusiones} recoge las conclusiones del
equipo.

\chapter{Desarrollo}
\label{cap:desarrollo}

\section{¿En qué consiste el paradigma?}

La computación neuromórfica diseña circuitos electrónicos que reproducen
principios de organización del tejido nervioso. Sus tres rasgos distintivos
frente a la arquitectura convencional son los siguientes.

\begin{itemize}
  \item \textbf{Colocalización de memoria y cómputo.} Las sinapsis artificiales
  almacenan los pesos en el mismo punto donde se realiza la operación, lo que
  elimina el tránsito constante de datos entre procesador y memoria.
  \item \textbf{Comunicación por impulsos.} La información viaja en forma de
  impulsos discretos, o \emph{spikes}, en lugar de valores continuos.
  \item \textbf{Operación dirigida por eventos.} No existe un reloj global que
  active todo el sistema en cada ciclo: una unidad consume energía únicamente
  cuando recibe o emite un impulso. De ahí procede la mayor parte del ahorro
  energético del paradigma \parencite{indiveri2015}.
\end{itemize}

\section{Fundamentos teóricos}

El punto de partida conceptual se encuentra en los sistemas VLSI analógicos
propuestos por \textcite{mead1990}, que reproducían en silicio la dinámica de las
neuronas biológicas aprovechando el comportamiento físico de los transistores en
régimen de subumbral.

El modelo de cómputo asociado son las redes neuronales de impulsos. En ellas, cada
neurona acumula las entradas que recibe a lo largo del tiempo y emite un impulso
cuando su potencial interno supera un umbral, tras lo cual el potencial se
reinicia. El modelo más utilizado es el de integración y disparo con fugas, cuya
dinámica puede expresarse como

\begin{equation}
\tau_m \frac{dV(t)}{dt} = -\bigl(V(t) - V_{\mathrm{reposo}}\bigr) + R\, I(t),
\label{eq:lif}
\end{equation}

donde $V(t)$ es el potencial de membrana, $\tau_m$ la constante de tiempo, $R$ la
resistencia de membrana e $I(t)$ la corriente de entrada. Cuando $V(t)$ alcanza
el umbral $V_{\mathrm{umbral}}$ se genera un impulso y el potencial vuelve a su
valor de reposo.

\begin{figure}[h!]
\centering
\begin{tikzpicture}[font=\small,scale=0.95]
  \draw[->,gray] (0,0) -- (9.4,0) node[right,color=gray]{$t$};
  \draw[->,gray] (0,-0.1) -- (0,2.3) node[above,color=gray]{$V$};
  \draw[dashed,red!60!black] (0,1.7) -- (9.2,1.7);
  \node[color=red!60!black,anchor=south west,font=\footnotesize] at (0.1,1.74) {umbral};
  \draw[azul,line width=1.1pt]
    (0,0.25) -- (1.1,0.85) -- (1.15,0.7) -- (2.2,1.4) -- (2.25,1.25)
    -- (3.1,1.71) -- (3.12,0.25) -- (4.3,0.8) -- (4.35,0.65)
    -- (5.4,1.2) -- (5.45,1.05) -- (6.6,1.71) -- (6.62,0.25)
    -- (8.0,0.75) -- (8.05,0.63) -- (9.0,1.1);
  \draw[line width=1.3pt,azul] (3.12,1.95) -- (3.12,2.4);
  \draw[line width=1.3pt,azul] (6.62,1.95) -- (6.62,2.4);
  \node[color=azul,anchor=west,font=\footnotesize] at (6.85,2.25) {impulsos emitidos};
\end{tikzpicture}
\caption{Potencial de membrana a lo largo del tiempo. Al alcanzar el umbral se
emite un impulso y el potencial se reinicia.}
\label{fig:neurona}
\end{figure}

La diferencia esencial con las redes neuronales artificiales convencionales es que
aquí el tiempo forma parte del cómputo: la información no se codifica solo en la
intensidad de una señal, sino en el instante y la frecuencia con que se producen
los impulsos.

\section{Aplicaciones prácticas}

El paradigma ha dado lugar a varias plataformas de hardware con planteamientos
distintos:

\begin{itemize}
  \item \textbf{TrueNorth} (IBM), un circuito digital de un millón de neuronas
  programables orientado a inferencia de muy bajo consumo \parencite{merolla2014}.
  \item \textbf{Loihi} (Intel), que incorpora aprendizaje en el propio chip
  mediante reglas de plasticidad sináptica \parencite{davies2018}.
  \item \textbf{SpiNNaker} (Universidad de Mánchester), una arquitectura masivamente
  paralela concebida para simular en tiempo real redes neuronales a gran escala
  \parencite{furber2014}.
\end{itemize}

Los ámbitos de aplicación más explorados son el reconocimiento de patrones
sensoriales con restricciones de energía (visión y audición artificial), la
robótica autónoma, el procesamiento en el borde de la red y las interfaces y
prótesis neuronales \parencite{schuman2022}.

\section{Cuatro dominios con resultados medidos}

La lista anterior de ámbitos resulta insuficiente para valorar la madurez real
del paradigma. A continuación se examinan cuatro dominios concretos, cada uno
con el resultado publicado que lo respalda.

\subsection{Navegación autónoma por visión}

Es el dominio con mejores resultados comparativos. En tareas de estimación de
flujo óptico sobre vehículos y drones, las redes de impulsos alcanzan un
\textbf{20\% menos de error} que redes convencionales de tamaño equivalente,
con unas \textbf{diez veces menos de energía} y hasta 48 veces menos parámetros
\parencite{chowdhury2025}. El dato relevante es doble: no solo ahorra energía,
también mejora la precisión.

\subsection{Sensores siempre activos}

En detección de palabra clave en audio sin conexión a la red, el procesador
Xylo Audio 2 consume \textbf{219 microvatios en reposo} y 0,0066 milijulios por
inferencia, con un acierto del 95,3\%. Esa cifra supera en aproximadamente un
orden de magnitud a Loihi y SpiNNaker en energía por inferencia
\parencite{bos2024}.

\subsection{Inteligencia artificial en centro de datos}

Hala Point, el mayor sistema neuromórfico construido hasta la fecha, integra
1.152 chips Loihi 2 con 1.150 millones de neuronas y 128.000 millones de
sinapsis, con un consumo máximo de \textbf{2.600 vatios} y una eficiencia
superior a 15 TOPS/W \parencite{intel2024}.

\subsection{Neuroprótesis}

Es el dominio menos maduro, y conviene afirmarlo con claridad.
\textcite{donati2024} concluyen que las tecnologías neuromórficas
\enquote{todavía tienen un largo camino por recorrer} y que no existen
implementaciones clínicas consolidadas: la revisión identifica líneas de
investigación abiertas, no sistemas desplegados.

\section{Lectura crítica de las cifras}

La tabla \ref{tab:cifras} reúne los cuatro resultados anteriores. Su
elaboración es aportación propia del equipo a partir de las cuatro fuentes
citadas.

\begin{table}[h!]
\centering
\caption{Resultados medidos por dominio de aplicación.}
\label{tab:cifras}
\begin{tabularx}{\textwidth}{L{0.8} L{0.7} L{1.5}}
\toprule
\textbf{Dominio} & \textbf{Plataforma} & \textbf{Dato medido} \\
\midrule
Visión autónoma & SNN frente a ANN & $-20$\% error; unas 10 veces menos energía \\
Audio siempre activo & Xylo Audio 2 & 219 $\mu$W en reposo; 0,0066 mJ por inferencia \\
Centro de datos & Hala Point & $1{,}15 \cdot 10^{9}$ neuronas en 2.600 W \\
Neuroprótesis & --- & sin cifras clínicas publicadas \\
\bottomrule
\end{tabularx}
\end{table}

Estas cifras exigen tres advertencias.

En primer lugar, existe un \textbf{sesgo de selección}. Ninguna de las fuentes
compara el paradigma en una tarea donde salga perdiendo. Los datos son
correctos, pero no representan al conjunto de aplicaciones posibles.

En segundo lugar, importa el \textbf{origen de la fuente}. La afirmación de que
los sistemas basados en Loihi operan con cien veces menos energía que las
arquitecturas convencionales procede de una nota de prensa del fabricante
\parencite{intel2024}, no de una publicación revisada por pares. Debe citarse
como tal. Las otras tres cifras sí proceden de revistas con revisión.

En tercer lugar, \textbf{no existe un banco de pruebas común}. Que Xylo supere
a Loihi y a SpiNNaker en un orden de magnitud en energía por inferencia no
significa que sea superior: significa que resuelven problemas distintos y que
no hay una medida estándar que permita compararlos. Esta carencia enlaza
directamente con el primer desafío tecnológico que se expone en el capítulo
siguiente.

\equipo{Buscar un estudio independiente que contradiga o matice alguna de
estas cuatro cifras. Encontrarlo reforzaría mucho esta sección; no encontrarlo
también es un hallazgo y conviene decirlo.}

\chapter{Resultados}
\label{cap:resultados}

Este capítulo recoge el análisis del equipo en sus dos vertientes: los desafíos
técnicos que hoy frenan la adopción del paradigma y las implicaciones éticas que
se derivan de sus características propias.

\section{Desafíos tecnológicos}

\subsection{Ausencia de un modelo de programación maduro}

La computación convencional dispone de un modelo de referencia estable y de
décadas de herramientas construidas sobre él. La computación neuromórfica carece
de un equivalente consolidado: cada plataforma de hardware impone su propio
entorno, sus primitivas y sus limitaciones, de modo que el código escrito para una
no es trasladable a otra. Esta fragmentación eleva el coste de entrada y dificulta
comparar resultados entre grupos de investigación \parencite{schuman2022}.

\subsection{Entrenamiento de las redes de impulsos}

El algoritmo de retropropagación, base del aprendizaje profundo, exige que la
función de activación sea derivable. El impulso no lo es: se trata de un evento
discreto. En consecuencia, el entrenamiento recurre a aproximaciones, como
sustituir la derivada por una función sustituta, convertir redes convencionales ya
entrenadas o emplear reglas de plasticidad local inspiradas en la biología. Ninguna
de estas vías alcanza todavía, de forma general, la precisión de los modelos
convencionales en tareas equivalentes \parencite{roy2019}.

\subsection{Fabricación y variabilidad de los dispositivos}

Las implementaciones analógicas y las basadas en memristores presentan dispersión
entre unidades fabricadas en el mismo proceso, además de sensibilidad a la
temperatura y al envejecimiento. Dos chips nominalmente idénticos pueden producir
resultados distintos, lo que compromete la reproducibilidad y obliga a introducir
mecanismos de calibración o de tolerancia a la variación
\parencite{indiveri2015}.

\section{Implicaciones éticas}

\subsection{Opacidad y rendición de cuentas}

Un sistema que aprende de forma continua, distribuida y local modifica su
comportamiento después del despliegue, sin un conjunto de parámetros fijo que
pueda inspeccionarse. Esto dificulta reconstruir por qué se tomó una decisión
concreta y, con ello, atribuir responsabilidad cuando el sistema falla. La
dificultad es mayor que en los modelos convencionales, donde al menos los pesos
permanecen estables tras el entrenamiento.

\subsection{Vigilancia ubicua}

El bajo consumo es la principal virtud técnica del paradigma y, a la vez, el
origen de su riesgo social más inmediato: permite instalar sensores capaces de
analizar imagen y sonido de forma permanente, alimentados por baterías de larga
duración o por captación de energía ambiental. La viabilidad económica de una
vigilancia continua en el espacio público deja de ser una barrera técnica para
convertirse en una decisión política.

\subsection{La frontera con lo humano}

Al describirse mediante vocabulario neurológico y al aplicarse a interfaces y
prótesis que se conectan con el sistema nervioso, el paradigma reabre preguntas
sobre la autonomía personal, la identidad y la privacidad mental.
\textcite{yuste2017} plantean en este sentido prioridades éticas específicas para
las neurotecnologías, entre ellas la protección de la identidad y de la agencia
individual.

\equipo{Decidir si se añade una cuarta implicación ética o si se
profundiza en alguna de las tres. Si se añade, buscar una fuente que la
respalde.}

\chapter{Discusión}
\label{cap:discusion}

\section{Estrategias para abordar los riesgos éticos}

Las tres medidas que proponemos se corresponden, una a una, con las implicaciones
expuestas en el capítulo anterior.

\subsection{Auditabilidad por diseño}

Frente a la opacidad, la medida no consiste en renunciar al aprendizaje continuo,
sino en exigir que todo sistema desplegado en un ámbito sensible registre de forma
verificable las entradas recibidas, los cambios aplicados a sus parámetros y las
decisiones emitidas. Ese registro debe estar disponible para auditores externos y
conservarse durante un plazo definido. La condición ha de incorporarse en la fase
de diseño, porque añadirla después resulta técnicamente inviable.

\subsection{Regulación del despliegue en espacios públicos}

Frente a la vigilancia ubicua, la respuesta corresponde al ámbito normativo:
delimitar en qué lugares y con qué finalidades puede instalarse un sensor con
capacidad de análisis permanente, establecer la obligación de señalizarlo y
prohibir el reconocimiento biométrico indiscriminado. Marcos como el Reglamento
europeo de inteligencia artificial ofrecen un punto de partida que puede adaptarse
a otros contextos regulatorios.

\subsection{Gobernanza específica de la neurotecnología}

Frente a los riesgos asociados a las interfaces neuronales, proponemos exigir
revisión por comités éticos independientes, consentimiento informado que
contemple de manera explícita el uso posterior de los datos neuronales, y un
tratamiento jurídico diferenciado de esos datos, que no pueden equipararse a
cualquier otro dato personal \parencite{yuste2017}.

\section{Impacto en la sociedad y en la humanidad}

\subsection{Dimensión ambiental y económica}

El argumento central a favor del paradigma es energético. Si una fracción
significativa de las tareas de inferencia migrara a hardware dirigido por
eventos, la presión sobre los centros de datos disminuiría de forma apreciable.
Conviene, no obstante, tratar las cifras con cautela: las comparaciones
publicadas suelen referirse a cargas de trabajo concretas y favorables, y no son
extrapolables al conjunto de aplicaciones \parencite{schuman2022}.

\equipo{Si se citan cifras de consumo de centros de datos, verificarlas
en la fuente original y decir a qué carga de trabajo corresponden. Sin ese
dato, la cifra no significa nada.}

\subsection{Dimensión social y de acceso}

El procesamiento autónomo en el propio dispositivo abre posibilidades claras en
prótesis, asistencia a personas con discapacidad y dispositivos médicos capaces
de operar sin conexión permanente. Al mismo tiempo, mientras la fabricación siga
concentrada en pocos actores, el acceso a esta tecnología reproducirá las
asimetrías ya existentes en la industria de los semiconductores.

\equipo{Aterrizar el impacto a México o Latinoamérica: qué significaría aquí
en concreto, con algún dato de la región si se encuentra.}

\subsection{Dimensión filosófica}

Imitar la arquitectura del cerebro no equivale a reproducir lo que el cerebro
hace. Conviene mantener esa distinción, porque el vocabulario del paradigma
(neuronas, sinapsis, aprendizaje) induce con facilidad a atribuir a estos sistemas
capacidades de las que carecen. El interés filosófico del paradigma no reside en
que nos acerque a construir una mente artificial, sino en que obliga a precisar
qué entendemos por procesar información y en qué se distingue de comprenderla.

\chapter{Conclusiones}
\label{cap:conclusiones}

El análisis realizado permite al equipo formular cuatro conclusiones.

\textbf{Primera.} La computación neuromórfica no consiste en acelerar el modelo
convencional, sino en reorganizarlo. Al unir memoria y cómputo y operar por
eventos, elimina el cuello de botella que limita la arquitectura de von Neumann,
a costa de renunciar a la generalidad de esta última.

\textbf{Segunda.} Sus obstáculos principales ya no son de fabricación, sino de
programación y entrenamiento. Existe hardware funcional, pero falta un modelo de
programación común y un método de aprendizaje que iguale, de forma general, a los
modelos convencionales.

\textbf{Tercera.} Las implicaciones éticas no son un añadido externo al
paradigma: proceden directamente de su mayor virtud técnica. Un consumo
energético tan reducido es, a la vez, lo que permite desplegar sensores
permanentes en cualquier lugar.

\textbf{Cuarta.} Las medidas de mitigación deben incorporarse durante el diseño.
La auditabilidad no puede añadirse a un sistema que ya aprende de forma continua,
y la regulación del despliegue pierde eficacia una vez que la infraestructura
está instalada.

\equipo{Acordar entre los cuatro la postura del equipo: ¿sustituirá al
modelo convencional, o será un acelerador especializado que conviva con él?
Defenderla con lo expuesto en los capítulos anteriores, no con opiniones.}

    %--- Referencias (formato APA, al final del documento) ---%
    \chapter*{Referencias}
    \addcontentsline{toc}{chapter}{Referencias}
    % Interlineado sencillo para que las referencias quepan en una hoja.
    \begin{spacing}{1}
        \printbibliography[heading=none]
    \end{spacing}

\end{document}
